\scanpage{184}
et par suite un isom canonique
\[
 \det(Q)^{\otimes2}\simeq{\Omega^1_{X^{p/S}/S}}^{\otimes(p-1)}\otimes\Lambda
 \tag*{($\ast\ast$)}
\]

\uline{Cas} 1 \quad $p=2$, on trouve
\[
 \deleted{\det Q^{\otimes}\simeq}\ \Omega^1_{X^{p/S}/S}\otimes\Lambda\simeq\det(Q)^{\otimes2}
\]
donc $\Omega^1_{X^{p/S}/S}\otimes\Lambda$ est \uline{canoniquement un carré}. Si $S$ est parfait, \inserted{et $X$ réduit, i.e. $\Lambda\simeq\underline O_{X^{p/S}}$,} ceci implique que $\Omega^1_{X/S}$ est lui-m\textsuperscript{e} canoniquement un carré. \deleted{Je doute que \unsure{ce|ceci} pour $X$} \deleted{\unsure{affine}} \deleted{(cf. résultat de)} Ceci n'est pas vrai si $S=\operatorname{Spec}k$, $k$ \inserted{corps} imparfait de car. 2, \inserted{m\textsuperscript{e} si $X$ lisse sur $k$,} comme \deleted{\unsure{et}} on voit \deleted{\unsure{déjà}} par l'exemple de Serre-Tate-\unsure{Fröhlich} : il suffit de prendre l'exemple d'une courbe \inserted{lisse propre} rationnelle sans point -- \margin{?} \deleted{\illegible{1}} il y en a sur tout corps imparfait de car 2

\uline{Cas} 2 \quad $p\ne2$, donc $\frac{p-1}{2}$ est un entier. On aimerait alors préciser $(\ast\ast)$ en un isom can
\[
 \psi:\det(Q)\simeq{\Omega^1_{X^{p/S}/S}}^{\otimes(\frac{p-1}{2})}\otimes\Lambda
 \left[\simeq{\Omega^1_{X/S}}^{\otimes(\frac{p-1}{2})}\otimes_{\underline O_X}\Lambda\right]
\]
[Sinon, il y aurait un \unsure{revêtement} de groupe $\mu_2\simeq\mathbb Z/2\mathbb Z$ de $X^{p/S}$, canoniquement associé à la situation]. Reprenons la base de $Q$, $(\dot x^i)_{1\le i\le p-1}$, \deleted{\unsure{et soit|et comme}} d'où une base de \deleted{\unsure{la}} $\det Q$, $\dot x\wedge\dot x^2\wedge-\wedge\dot x^{p-1}$, et la valeur de son carré, par l'isom $(\ast\ast)$, est $\Delta(d\widetilde x^{\otimes(p-1)}\otimes\Lambda)$, où $\Delta$ est le déterminant de la matrice écrite plus haut, savoir $(-1)^{(p-1)-1}1.2.\ldots(p-1)=\deleted{\unsure{+}}-(1.2.\ldots(p-1))$. Comme $\mathbb F_p^\ast$ est un groupe cyclique d'ordre \deleted{\unsure{un}} pair \unsure{pour|$p-1$} le produit de ses éléments est $-1$, donc on trouve $\Delta=\unsure{+1|1}$
% archive sheet number: 183
